\documentclass[10pt,a4paper]{amsart}
\usepackage[margin=19mm]{geometry}
\usepackage{amsmath,amssymb,mathtools}
\usepackage[scr=rsfs,cal=euler]{mathalfa}
\usepackage{xfrac}
\usepackage[unicode,hidelinks,bookmarks=false]{hyperref}
\newcommand{\ie}{i.e.\ }
\newcommand{\cf}{cf.\ }
\newcommand{\resp}{respectively\ }
\newcommand{\Subsection}{\subsection}
\providecommand{\nobreakdash}{\nobreak\hbox{-}}
\setlength{\emergencystretch}{2em}
\multlinegap0pt
\usepackage{bm}
\usepackage{bbm}
\usepackage{dsfont}
\usepackage{xspace}
\usepackage{cancel}
\usepackage{mathtools}
\usepackage{amsthm}
\makeatletter
\@ifundefined{lemma}{\newtheorem{lemma}{Lemma}}{}
\@ifundefined{proposition}{\newtheorem{proposition}[lemma]{Proposition}}{}
\@ifundefined{theorem}{\newtheorem{theorem}[lemma]{Theorem}}{}
\makeatother
\title[On Generalized Rickart $*$-rings]{On Generalized Rickart $*$-rings}
\author{Anil Khairnar, Sanjay More}
\date{Source: arXiv:2508.20842v1 (CC BY 4.0)}
\begin{document}
\maketitle

\noindent\emph{Source and licence.} On Generalized Rickart $*$-rings. Version arXiv:2508.20842v1 (CC BY 4.0), distributed under the Creative Commons Attribution 4.0 International licence, as recorded in the arXiv licence metadata for this exact version and shown on the abstract page https://arxiv.org/abs/2508.20842v1. Full rights notice: licenses/arxiv-2508.20842-CC-BY-4.0.txt.\par\smallskip

\noindent\textit{Editorial scope.} Counted unit: the Theorem whose source label is \texttt{None} in the article (source0.tex lines 501--508, source label \texttt{None}), delivered together with the article's own complete original proof of it (source0.tex lines 509--552, one \texttt{proof} environment of 44 lines). No mathematics is written, repaired or re-derived: statement and proof are the article's own text line for line. Required context retained uncounted: s2pr8 (lines 291--296). Every definition, display and label of those lines is retained unchanged. Omitted from this package and not redistributed: 1--290, 297--500, 553--797. Nothing inside a retained excerpt is omitted or altered: every retained body is delivered exactly as the article's lines are written, so no omission receipt of any kind accompanies this package. Citations: the retained excerpts carry no citation command at all, so no bibliography entry of the article is transcribed and no reference is redistributed. Reference adaptation: none was needed and none was made; every reference inside the delivered unit resolves inside it, so no unresolved reference or citation is produced. Printed slips kept literal: no printed word, symbol, hypothesis or number of the counted statement or of its proof was altered. Layout and class: the article's own class is \texttt{amsart} with packages including amssymb, latexsym, xy, eucal, mathrsfs, MnSymbol; the wrapper is the lane's pinned \texttt{amsart} editorial wrapper at 10pt on A4, which restates the theorem-like environments the retained text needs and adds the article's own macro definitions that the retained text needs (none), with extra wrapper packages (bm, bbm, dsfont, xspace, cancel, mathtools). The delivered numbering is the unit's own. Assets: no retained span contains a figure environment, a graphics-inclusion command, a PDF-page inclusion, a table or a TikZ picture; no image file is copied into this package, the package declares no asset dependency and no image file is redistributed with it. Licence: version arXiv:2508.20842v1 of this article is distributed under the Creative Commons Attribution 4.0 International licence, as recorded in the arXiv licence metadata for this exact version and shown on the abstract page https://arxiv.org/abs/2508.20842v1. A full rights notice is delivered at licenses/arxiv-2508.20842-CC-BY-4.0.txt. Characters: every retained excerpt is pure ASCII, so the delivered engine, XeLaTeX with its default Latin Modern text font, reports no missing character.
\begin{proposition} \label{s2pr8} The following conditions on a $*$-ring $R$ are equivalent.
	\begin{enumerate}
		\item $R$ is a generalized  Rickart $*$-ring.  
		\item $R$ is a generalized weakly Rickart $*$-ring with unity.
	\end{enumerate}
\end{proposition}

\begin{theorem}  A generalized weakly Rickart $*$-ring can be embedded in a generalized Rickart $*$-ring
provided there exists a ring $K$ such that
\begin{enumerate}
	\item  $K$ is an integral domain with involution.
	\item  $R$ is $*$-algebra over $K$. 
	\item  For any nonzero $\lambda\in K$, there exists a projection $e_\lambda$ such that $\lambda x=0$ implies   $GRP(x)\leq e_\lambda$.
\end{enumerate}
\end{theorem}

\begin{proof} Let $ R_1 = R\oplus K$ (the additive group direct sum) with operations as defined above. First we prove that for any self-adjoint element $a\in R$ and $0 \ne \lambda \in K$ there exists largest projection $g$ such that $(ag+\lambda g)^m=0$ for some $m\in\mathbb N $. Let $GRP(a)=e_0$.
	Then  $a^me_0=a^m$ and $a^my=0$ implies $e_0y=0$ for some $m\in\mathbb N$. 
	Let $e_\lambda$ be a projection which exists by the assumption (3).
	Let $e$ be the largest projection in $\{e_0,e_\lambda\}$.
	Let $GRP(ae+\lambda e)=h$. Hence there exists $ m\in\mathbb N$ such that $(ae+\lambda e)^mh=(ae+\lambda e)^m$ and  $(ae+\lambda e)^my=0$ implies $hy=0$.  Now $e_0\leq e$ implies $e_0=e_0e=ee_0$.
	Let $g = e-h$.	Since $(ae+\lambda e)^me=(ae+\lambda e)^m$, we have $h\leq e$, that is $h=he =eh$. Therefore $eg=e-eh=e-h=g$. This gives $g\leq e$ and hence $g=eg=ge$.
	Thus $(ag+\lambda g)^m=(a eg+\lambda eg)^m=(ae+\lambda e)^mg
	=(ae+\lambda e)^m(e-h)=(ae+\lambda e)^m-(ae+\lambda e)^mh=0$.
	To prove $g$ is largest. Suppose $(ak+\lambda k)^m=0$. We have $e_0=e_0e=ee_0$. Since $a^me_0=a^m$, $a^me_0e=a^me$, which implies $a^me_0=a^me$. Therefore $ea^m=a^m$, which gives $kea^mk=ka^mk$. Hence $(ke-k)a^mk=0$. Since $(ak+\lambda k)^m=0$, we have $(ke-k)\{-\binom{m}{1}a^{m-1}k\lambda k-\cdots-\lambda^mk\}=0$. Equating coefficient of $a^{m-1}k$, we get  $\lambda m (ke-k)=0$. 
	Therefore $\lambda (ke-k)=0$. Let $GRP(ke-k)=f$. Then $(ke-k)^nf=	(ke-k)^n $ and $(ke-k)^ny=0$ implies  $fy=0$. Since $(ke-k)^ne=0$, we have $fe=0$. Further $\lambda (ke-k)=0$ implies $\lambda (ke-k)^n=0$.  Which gives $\lambda (ke-k)^nf=0$. Hence $ (ke-k)^n\lambda f=0$ . Therefore $f(\lambda f)=0$ . Thus $\lambda f=0$. By (3) $GRP(f)\leq e_\lambda\leq e$. Therefore $f\leq e$, that is $f=fe=ef$ (since $GRP(f)=f$). Hence $f=0$. As $(ke-k)^n= (ke-k)^nf$ implies $(ke-k)^n=0$. But $(ke-k)^n=\pm (ke-k)$. Hence $\pm (ke-k) =0$ gives $ke=k$. So $(ae+\lambda e)^m k=(aek+\lambda ek)^m=(ak+\lambda k)^m=0$. Therefore $hk=0$. Hence $kg=k(e-h)=ke-kh=k-0=k$. Thus $k\leq g$. Hence $g$ is the largest projection  such that $(ag+\lambda g)^m=0$.  Since $(0,1)$ is the unity element of $R_1$. By above Proposition \ref{s2pr8}  it is enough to show that $R_1$ is a generalized weakly Rickart $*$-ring. Let $(a,\lambda)\in R_1$.\\
	Suppose $\lambda=0$. Since $(a,0)\in R_1$, we have $a\in R$. 
	As $R$ is a generalized weakly Rickart $*$-ring, $GRP(a)=e$ exists in $R$.
	That is for some $n\in \mathbb N$, $a^ne=a^n$; and for $y \in R$, $a^ny=0$ implies $ey=0$.
	We will prove that $GRP(a,0)=(e,0)$. 
	Since $(a,0)^n(e,0)=(a^n,0)(e,0)=(a^ne,0)=(a^n,0)=(a,0)^n$.
	Suppose $(a,0)^n(b,\mu)=0$. Then $(a^n,0)(b,\mu)=0$.
	This implies $(a^nb+\mu a^n,0)=0$, that is $ a^nb+\mu a^n=0$.
	This gives $a^nb+\mu a^ne=0$, and hence $a^n (b+\mu e)=0$.
	Since  $a^ny=0$ implies  $ey=0$, we have $e(b+\mu e)=0$.
	Therefore $eb+\mu e=0$. That is $(e,0)(b,\mu)=0$. Hence $GRP(a,0)=(e,0)$.\\
	Now, suppose $\lambda\ne 0$. Then there exists a largest projection $g$ in $R$ such that $(ag+\lambda g)^t=0$ for some $t\in\mathbb N$. Note that $(-g,1)$ is a projection in $R_1$.
	We prove that $GRP(a,\lambda)=(-g,1)$.
	Consider $(a,\lambda)^t(-g,1)=(a^t+\binom {t}{1}a^{t-1}\lambda+\cdots+\binom {t}{t-1}a\lambda ^{t-1},\lambda^t)(-g,1) =(-a^tg-\binom {t}{1}a^{t-1}\lambda g-\cdots -\binom {t}{t-1}a\lambda^{n-1}g-\lambda^tg+a^t+\binom {t}{1}a^{t-1}\lambda+\cdots+\binom {t}{t-1}a\lambda^{t-1},\lambda^t)=(-(ag+\lambda g)^t+a^t+\binom {t}{1}a^{t-1}\lambda+\cdots+\binom {t}{t-1}a\lambda^{t-1},\lambda^t)=(a^t+\binom {t}{1}a^{t-1}\lambda+\cdots+\binom {t}{t-1}a\lambda^{t-1},\lambda^t)=(a,\lambda)^t$.
	To prove $(a,\lambda)^t(b,\mu)=0$ implies $(-g,1)(b,\mu)=0$. 
	Let $(a,\lambda)^t(b,\mu)=0$. Then $(a^t+\binom {t}{1}a^{t-1}\lambda+\cdots+\binom {t}{t-1}a\lambda^{t-1},\lambda^t)(b,\mu)=0$. 
	This implies $ \lambda^t\mu=0\Rightarrow\mu=0$  (since $\lambda\ne 0$). 
	Therefore  $(a^t+\binom {t}{1}a^{t-1}\lambda+\cdots+\binom {t}{t-1}a\lambda^{t-1},\lambda^t)(b,0)=0$.
	
	 To prove  $(a,\lambda)^t(b,0)^m=0$ implies $(-g ,1)(b,0)^m=0$ for some $m \in \mathbb N$.
	 Let $GLP(b)=f$, then there exist $m \in \mathbb N$ such that 
	 $fb^m=b^m$; and for $y \in R$, $yb^m=0$ implies $yf=0$. 
	Therefore $(a,\lambda)^t(b,0)^m=0$. This implies that 
	$(a^t+\binom {t}{1}a^{t-1}\lambda+\cdots+\binom {t}{t-1}a\lambda^{t-1},\lambda^t)(b^m,0)=0$.
	Therefore $(\{a^t+\binom {t}{1}\lambda a^{t-1}+\cdots+\binom {t}{t-1}\lambda^{t-1}a\}b^m+\lambda^tb^m,0)=0$.
	Hence $(\{a^tf+\binom {t}{1}\lambda a^{t-1}f+\cdots+\binom {t}{t-1}\lambda^{t-1}af+\lambda^tf\}b^m,0)=0$.
	This gives $(af+\lambda f)^tb^m=0$. 
	Since $yb^m=0$ implies $yf=0$, we have $(af+\lambda f)^tf=0$.
	This implies $(af+\lambda f)^t=0$.
	But $g$ is the largest projection such that $(ag+\lambda g)^t=0$.
	Therefore $f\leq g$, which gives  $f(1-g)=0$.
	Hence $(-g,1)(b,0)^m=(-g,1)(b^m,0)=(-gb^m+b^m,0)=((1-g)b^m,0)=((1-g)fb^m,0)=(0,0)=0$.
	
	Hence $GRP(a,\lambda)=(-g,1)$. Thus  $R_1$ is a  generalized Rickart $*$-ring.	
\end{proof}

\end{document}
